\documentclass[aps,prl,onecolumn,superscriptaddress,nofootinbib,longbibliography]{revtex4-2}

\usepackage{amsmath,amssymb,bm}
\usepackage{booktabs}
\usepackage{tabularx,array}
\usepackage{graphicx}
\usepackage{xcolor}
\usepackage{xurl}
\usepackage{hyperref}
\hypersetup{
  colorlinks=true,
  allcolors=blue!55!black,
  pdftitle={Supplemental Material: Basis-Dependent Sparsity and Energy Recovery in [2Fe--2S]},
  pdfauthor={Anonymous Authors},
  pdfsubject={Simulation protocol, numerical validation, and reproducibility data},
  pdfkeywords={iron-sulfur cluster, Davidson, configuration interaction, natural orbitals, reproducibility}
}

\newcommand{\Eh}{E_{\mathrm h}}
\newcommand{\cas}{\mathrm{CAS}(30e,20o)}
% Frozen numerical values generated by analysis/build_results_tex.py.
\newcommand{\FinalEnergy}{-116.60560912042631}
\newcommand{\FinalResidual}{6.9016\times10^{-7}}
\newcommand{\BenchmarkKFiveTwelveCoupling}{1.2105313}
\newcommand{\BenchmarkKFiveTwelveTail}{-0.6052656}
\newcommand{\NaturalKFiveTwelveRelaxedEnergy}{-116.486260278799}
\newcommand{\NaturalKFiveTwelveRelaxedErrorMilli}{119.349}
\newcommand{\NaturalKFiveTwelveRelaxationGainMilli}{49.634}
\newcommand{\NaturalKFiveTwelveProjectedResidual}{1.68\times10^{-13}}
\newcommand{\NaturalKFiveTwelveFullResidual}{0.311254}
\newcommand{\NaturalKFiveTwelveIndependentDifference}{5.68\times10^{-14}}
\newcommand{\NaturalReplayTrace}{29.999999987374895}
\newcommand{\NaturalReplayExpectedTrace}{30.000000000000426}
\newcommand{\NaturalReplayTraceError}{1.26\times10^{-8}}
\newcommand{\NaturalReplayRDMSymmetry}{0}
\newcommand{\NaturalReplayRDMOffDiagonal}{2.27\times10^{-11}}
\newcommand{\NaturalReplayRDMReferenceDifference}{6.27\times10^{-10}}
\newcommand{\NaturalReplayOccupationsDifference}{6.27\times10^{-10}}
\newcommand{\LateBasisSentence}{At \(K=4{,}194{,}304\), the natural-orbital projection increases \(W_K\) from 0.985678 to 0.999776 and reduces \(\Delta E_K\) from 17.705 to \(0.3488\,\mathrm{m}\Eh\); its residual nevertheless remains \(2.72\times10^{-2}\,\Eh\).}


\begin{document}

\title{Supplemental Material for ``Basis-Dependent Sparsity and Energy Recovery in a Strongly Correlated [2Fe--2S] Active Space''}
\author{Anonymous Authors}
\affiliation{Affiliations withheld during scientific review}
\maketitle

This Supplemental Material defines the active-space Hamiltonian, gives the
complete staged Davidson simulation, documents direct and density-matrix
replay, derives and validates the natural-orbital transformation, specifies the
top-\(K\) and Hamiltonian-decomposition experiments, and maps every result to
an executable script and frozen artifact.  ``Full configuration interaction''
below always means the complete determinant space of the specified \(\cas\)
Hamiltonian in the fixed \(M_S=0\) sector; it does not mean full-molecule or
complete-basis FCI.

\section{Hamiltonian, determinant convention, and public inputs}

The calculation uses a real, spin-independent FCIDUMP with 20 orthonormal
spatial orbitals, 30 electrons, \(M_S=0\), and \(E_{\mathrm{core}}=0\).  In the
chemist's integral convention implemented by PySCF,
\begin{equation}
 \hat H=\sum_{pq\sigma} h_{pq}a^{\dagger}_{p\sigma}a_{q\sigma}
 +\frac{1}{2}\sum_{pqrs\sigma\tau}(pq|rs)
 a^{\dagger}_{p\sigma}a^{\dagger}_{r\tau}a_{s\tau}a_{q\sigma}.
 \label{eq:s-hamiltonian}
\end{equation}
There are 15 alpha and 15 beta electrons.  Each spin sector contains
\(\binom{20}{15}=15{,}504\) occupation strings, so the CI tensor has shape
\(15{,}504\times15{,}504\), dimension 240,374,016, and occupies
1,922,992,256 bytes in little-endian float64 NPY format.

The primary FCIDUMP is retrieved from the immutable content associated with
Ref.~\cite{qatHamiltonian}; its SHA-256 is
\begin{quote}\small\ttfamily
bd6ef31773f7b082171f73faa6c98924948e2c250b6fbec74504d64361c8c6c7.
\end{quote}
The second, independently distributed QMB FCIDUMP is retrieved at repository
revision \path{ab1a7002ccbaca0ac43acafaf0f40400d14c3940} from
Ref.~\cite{qmbModels}; the uncompressed file has SHA-256
\begin{quote}\small\ttfamily
95d8786af06eeea2107e19ffd98c66a6ca97fc8c9864175a4f6d64512b6f2df9.
\end{quote}

\subsection{Orbital equivalence of the two FCIDUMPs}

The two files are not byte-identical because their active orbitals differ.
The script \path{code/compare_hamiltonian_orbital_equivalence.py} diagonalizes
each real symmetric one-electron matrix.  Its spectrum is nondegenerate, with
a minimum adjacent gap of \(1.6042\times10^{-2}\,\Eh\), so the ordered
eigenvectors are unique up to column signs.  After fixing the first sign and
determining the remaining signs from nonzero two-electron witnesses, the
orthogonal matrix \(U\) is tested against
\begin{align}
 h^{(B)}_{pq}&=\sum_{ij}U_{ip}h^{(A)}_{ij}U_{jq},
 \label{eq:s-rotation-h}\\
 g^{(B)}_{pqrs}&=\sum_{ijkl}U_{ip}U_{jq}U_{kr}U_{ls}g^{(A)}_{ijkl}.
 \label{eq:s-rotation-g}
\end{align}
The final comparison contracts the complete two-electron tensor without
screening.  Table~\ref{tab:equivalence} establishes orbital-rotation
equivalence to the numerical precision of the inputs; it does not attach an
uncertainty to a rounded literature energy.

\begin{table}[h]
\caption{Orbital-equivalence checks.  ``Maximum'' denotes the largest
elementwise absolute difference.}
\label{tab:equivalence}
\begin{ruledtabular}
\begin{tabular}{lc}
Quantity & Value \\
\hline
Maximum ordered \(h\)-eigenvalue difference & \(2.1316\times10^{-14}\,\Eh\) \\
Maximum defect of \(U^{\mathsf T}U=I\) & \(3.1086\times10^{-15}\) \\
Maximum rotated one-electron residual & \(3.5527\times10^{-14}\,\Eh\) \\
RMS complete two-electron-tensor residual & \(3.0805\times10^{-15}\,\Eh\) \\
Maximum complete two-electron-tensor residual & \(1.2055\times10^{-13}\,\Eh\) \\
\end{tabular}
\end{ruledtabular}
\end{table}

\section{Checkpointed Davidson simulation}

\subsection{Matrix-free Ritz iteration}

The \(240{,}374{,}016\)-dimensional Hamiltonian matrix is never formed.
PySCF's \texttt{fci.direct\_spin1} implementation applies the Slater--Condon
Hamiltonian to a CI tensor.  At update \(m\), an orthonormal trial basis
\(V_m\) defines
\begin{equation}
 T_m=V_m^{\mathsf T}HV_m.
\end{equation}
The lowest eigenpair \((\theta_m,y_m)\) of \(T_m\) gives the Ritz vector
\(c_m=V_my_m\) and residual \(r_m=Hc_m-\theta_mc_m\).  A diagonally
preconditioned correction expands the basis; the iteration restarts when the
configured subspace capacity is reached.  We set
\texttt{davidson\_only=True} and \texttt{pspace\_size=800} so the reported path
is iterative throughout.

The implementation in \path{code/continue_fci_checkpointed.py} adds three
reproducibility safeguards.  First, it hashes the FCIDUMP and input vector
before numerical work.  Second, every eighth callback writes the current Ritz
vector to a temporary NPY, flushes and synchronizes it, and atomically renames
it.  Third, after each stage it reloads the saved endpoint, constructs \(Hc\),
and recomputes the norm, Rayleigh quotient, and residual.  Consequently, every
reported endpoint is an observable of a concrete coefficient tensor rather
than an extrapolation or a printed scalar.

\subsection{Complete staged protocol}

Table~\ref{tab:stages} gives every stage used to obtain the reported state.
Stage A starts from PySCF's default determinant-space initial guess.  B1 tests
and establishes a wider 64-vector restart.  B2 creates the synchronized input
for the 160-update production stage C.  Only the explicit vector listed for a
stage is passed to the next one.

\begin{table*}[t]
\caption{Davidson stages.  Energies and residuals are direct contractions of
the saved endpoint.  The update limit is a numerical stop, not a convergence
claim.}
\label{tab:stages}
\small
\setlength{\tabcolsep}{3.2pt}
\begin{tabularx}{\textwidth}{@{}>{\raggedright\arraybackslash}Xrrcccl@{}}
\toprule
Stage & Updates & Space & Energy (\(\Eh\)) & Residual (\(\Eh\)) & \(E\) tol. & Vector SHA-256 prefix \\
\midrule
A: initialization & 110 & 14 & \(-116.6055095343891\) & \(2.8201\times10^{-3}\) & \(10^{-10}\) & \texttt{3b20ec2e1aa3} \\
B1: wide restart & 2 & 64 & \(-116.6055107620259\) & \(2.0814\times10^{-3}\) & \(10^{-13}\) & \texttt{cc9350d4bdad} \\
B2: checkpoint & 16 & 64 & \(-116.6055790186671\) & \(4.1575\times10^{-3}\) & \(10^{-13}\) & \texttt{1664c06a55c7} \\
C: production & 160 & 64 & \(-116.60560912042631\) & \(6.9016\times10^{-7}\) & \(10^{-13}\) & \texttt{45e63bb63cd1} \\
\bottomrule
\end{tabularx}
\end{table*}

Stages B and C used the residual target \(10^{-8}\,\Eh\), 64 software
threads, a 512-GiB node allocation, and reported maximum memory 450,000 MB.
The production eigensolver wall time was 23,387.508 s.  The complete driver,
including initial and final contractions, checkpoint input/output, and hashing,
took 23,494.156 s.  The final solver flag is
\texttt{converged=false}: the 160-update limit was exhausted while the residual
remained above its target.

The final coefficient tensor is \path{fci_ritz_vector.npy}, with SHA-256
\begin{quote}\small\ttfamily
45e63bb63cd10e953260b6f5dabb8a1dda5518fbf5959002c23888dacb599945.
\end{quote}
Every downstream script asserts this hash, the FCIDUMP hash, the tensor shape,
and float64 dtype before evaluating a reported quantity.

\section{Fresh-process replay and numerical stability}

The script \path{code/verify_publication_state.py} starts in a fresh Python
process and does not import the optimization driver.  It evaluates the energy
by two paths.  The first absorbs the one-electron integrals into the
two-electron contraction and explicitly constructs \(H_{\mathrm{el}}c\):
\begin{align}
 N&=\langle c|c\rangle,\\
 E_{Hc}&=E_{\mathrm{core}}+
 \frac{\langle c|H_{\mathrm{el}}|c\rangle}{N},\\
 \rho&=\frac{\lVert H_{\mathrm{el}}c-(E_{Hc}-E_{\mathrm{core}})c\rVert_2}
 {\sqrt N}.
\end{align}
The second constructs separate one- and two-particle reduced density matrices
and evaluates
\begin{equation}
 E_{\mathrm{RDM}}=E_{\mathrm{core}}
 +\sum_{pq}h_{pq}\gamma_{qp}
 +\frac{1}{2}\sum_{pqrs}g_{pqrs}\Gamma_{pqrs}.
 \label{eq:s-rdm}
\end{equation}
Both paths use PySCF, but they have distinct contraction routes; ``replay'' does
not mean a second electronic-structure library.  Results appear in
Table~\ref{tab:replay}.

\begin{table}[h]
\caption{Fresh-process observables of the saved state.}
\label{tab:replay}
\begin{ruledtabular}
\begin{tabular}{lc}
Quantity & Value \\
\hline
Squared norm & \(0.999999999999999\) \\
Direct Rayleigh energy & \(-116.60560912042631\,\Eh\) \\
RDM-reconstructed energy & \(-116.60560912042571\,\Eh\) \\
Absolute energy-path difference & \(5.9686\times10^{-13}\,\Eh\) \\
Residual norm & \(6.9016\times10^{-7}\,\Eh\) \\
Energy variance \(\rho^2\) & \(4.7631\times10^{-13}\,\Eh^2\) \\
\(\langle \hat S^2\rangle\) & \(2.8570\times10^{-9}\) \\
Spin multiplicity & \(1.0000000057\) \\
\end{tabular}
\end{ruledtabular}
\end{table}

Two PySCF kernels also accumulated the spin-summed one-particle RDM through
different paths.  On the 240-million-component tensor their maximum
elementwise difference was \(1.7360\times10^{-10}\), at diagonal element
\((16,16)\), and therefore exceeded an initially chosen \(10^{-10}\) absolute
comparison.  The subsequent stability calculation found a relative Frobenius
difference \(8.4093\times10^{-11}\), zero reported symmetry defect in either
RDM, and traces 29.999999997413795 and 29.999999999999968.  Equation
(\ref{eq:s-rdm}) retained sub-pico-Hartree agreement with the direct energy.
We report this as an accumulation-order stability study rather than replacing
the initial comparison by a looser undeclared threshold.

\section{Exact natural-orbital representation of the same state}

The real symmetric, spin-summed one-particle RDM is symmetrized and
diagonalized,
\begin{equation}
 \frac{\gamma+\gamma^{\mathsf T}}{2}U
 =U\,\mathrm{diag}(n_1,\ldots,n_{20}),
 \qquad U^{\mathsf T}U=I.
\end{equation}
The natural occupations sum to 29.9999999974.  Ten values are strongly
fractional, ranging from 0.604296 to 1.468175; the complete ordered list is
stored in \path{data/natural_occupations.csv}.  The maximum off-diagonal RDM
element after rotation is \(1.07\times10^{-15}\), and the maximum
orthogonality defect of \(U\) is \(2.44\times10^{-15}\).

The integral tensors are transformed with the same \(U\) according to
Eqs.~(\ref{eq:s-rotation-h}) and (\ref{eq:s-rotation-g}).  The determinant
transformation for one spin sector is
the exterior-power representation of \(U\): a matrix element between two
occupation strings is the corresponding 15-by-15 minor.  Direct evaluation of
all such minors is unnecessarily expensive because the sector contains only
five holes.  Jacobi's complementary-minor identity,
\begin{equation}
 \det U_{I,J}=(-1)^{\Sigma I+\Sigma J}\det(U)
 \det\!\left[(U^{-1})_{J^c,I^c}\right],
 \label{eq:jacobi}
\end{equation}
maps every 15-electron minor to a 5-hole minor.  The script
\path{analysis/transform_ci_natural_basis.py} maps the electron tensor to the
hole ordering, calls PySCF's determinant transformation in the five-hole
sector, restores electron ordering and parity, and writes the transformed
float64 tensor.  A seeded six-orbital, four-electron test compares the direct
electron and complementary-minor paths and gives maximum difference
\(2.50\times10^{-16}\).

The source and transformed tensors are contracted in the same hash-locked
process before any sparsity analysis.  Table~\ref{tab:natural-gates} shows that
the transformation changes representation but not the physical state.  The
natural-basis FCIDUMP, rotation, and transformed vector have SHA-256 values
\texttt{adb401872023...}, \texttt{3a3dcbdfac4c...}, and
\texttt{3855dcb7ce8c...}, respectively; their full hashes are recorded in the
manifest.

\begin{table}[h]
\caption{Same-process orbital-transformation checks.}
\label{tab:natural-gates}
\begin{ruledtabular}
\begin{tabular}{lccc}
Quantity & Benchmark basis & Natural basis & Difference \\
\hline
Squared norm & 0.999999999999999 & 0.999999999999985 & \(1.39\times10^{-14}\) \\
Energy (\(\Eh\)) & \(-116.60560912042631\) & \(-116.60560912042600\) & \(3.13\times10^{-13}\) \\
Residual (\(\Eh\)) & \(6.9015577054\times10^{-7}\) & \(6.9015577065\times10^{-7}\) & \(1.12\times10^{-16}\) \\
\end{tabular}
\end{ruledtabular}
\end{table}

An independent frozen-vector calculation further tests the direction of the
orbital transformation, rather than only its spectrum.  Before reading either
FCIDUMP or CI tensor, it verifies the SHA-256 of the benchmark-basis occupation
list, 
\texttt{2e9632ff525337eae6329eac5cd5e39b71bd08af85291a335033296b5a732b7d}.
It then constructs the one-particle RDM directly in the natural representation
and compares the complete symmetrized matrix elementwise with the diagonal
matrix made from that frozen list.  The final trace is
\(\NaturalReplayTrace\), compared with
\(N\langle c|c\rangle=\NaturalReplayExpectedTrace\); the absolute difference
is \(\NaturalReplayTraceError\).  The maximum raw symmetry defect,
off-diagonal element, complete-matrix difference, and occupation-eigenvalue
difference are \(\NaturalReplayRDMSymmetry\),
\(\NaturalReplayRDMOffDiagonal\),
\(\NaturalReplayRDMReferenceDifference\), and
\(\NaturalReplayOccupationsDifference\), respectively.

A preliminary replay used an absolute trace threshold of \(10^{-8}\) and
returned a trace of 29.999999987374895, approximately
\(1.26\times10^{-8}\) from \(N\langle c|c\rangle\).  Before the final replay,
we fixed a \(2\times10^{-8}\) absolute tolerance and added the stronger
direction-sensitive complete-matrix comparison above; neither threshold nor
test was altered after observing the final output.  This records the
accumulation-level numerical tolerance without converting it into a physical
uncertainty on the energy.

\section{Top-\texorpdfstring{\(K\)}{K} projections and exact decomposition}

For each orbital representation, the streaming ranking in
\path{analysis/analyze_state_compressibility.py} retains the exact global
largest 4,194,304 values of \(|c_I|^2\).  At each input block it keeps the
block's local candidates, merges them with the current global set, and applies
an exact partition; a final stable sort fixes descending order.  The 17 support
sizes are
\begin{multline}
 K=64,128,256,512,1024,2048,4096,8192,16384,32768,\\
 65536,131072,262144,524288,1048576,2097152,4194304.
\end{multline}
For every \(K\), a zeroed full-shaped tensor receives the selected
coefficients without reoptimization and is normalized.  A fresh \(H\phi_K\)
contraction gives its Rayleigh energy and full-space residual.  The program
requires nested cumulative norm and
\(E_K\ge E_\star-10^{-10}\,\Eh\), and exports all points as CSV.  This protocol
never converts norm loss into an assumed energy loss.

For the benchmark basis, the coefficient probabilities
\(p_I=|c_I|^2\) give
\begin{equation}
 N_{\mathrm{IPR}}=(\sum_Ip_I^2)^{-1}=5525.89,\qquad
 N_{\mathrm{Shannon}}=\exp[-\sum_Ip_I\ln p_I]=109228.69.
\end{equation}
These scalar counts are representation dependent.  The direct two-basis
comparison is summarized in Fig.~2 and the complete files
\path{data/topk_tracker_basis.csv} and
\path{data/topk_natural_basis.csv}.  \LateBasisSentence

\begin{table*}[h]
\caption{Representative amplitude-ranked projections.  Energy errors are
relative to the full Ritz state and every residual is evaluated in the complete
active-space determinant space.}
\label{tab:representative-topk}
\small
\begin{tabular}{rcccccc}
\toprule
& \multicolumn{3}{c}{Benchmark orbitals} & \multicolumn{3}{c}{Natural orbitals} \\
\cmidrule(lr){2-4}\cmidrule(lr){5-7}
\(K\) & \(W_K\) & \(\Delta E_K\) (m\(\Eh\)) & residual (\(\Eh\))
& \(W_K\) & \(\Delta E_K\) (m\(\Eh\)) & residual (\(\Eh\)) \\
\midrule
64 & 0.075498 & 784.161 & 0.840786 & 0.418118 & 406.615 & 0.688659 \\
512 & 0.191346 & 605.266 & 0.785570 & 0.768507 & 168.983 & 0.472888 \\
65,536 & 0.726934 & 212.378 & 0.546320 & 0.974209 & 25.545 & 0.207036 \\
4,194,304 & 0.985678 & 17.705 & 0.186547 & 0.999776 & 0.3488 & 0.027168 \\
\bottomrule
\end{tabular}
\end{table*}

To decompose the error invariantly under an additive Hamiltonian constant, set
\(A=H-E_\star I\), \(p=P_Kc\), \(q=(I-P_K)c\), and
\(W=\langle p|p\rangle\).  Because \(c=p+q\) and
\(\langle c|A|c\rangle=0\),
\begin{align}
 W(E_K-E_\star)&=\langle p|A|p\rangle\\
 &=-\langle q|A|q\rangle
 -2\operatorname{Re}\langle p|A|q\rangle.
 \label{eq:s-decomposition}
\end{align}
The analysis contracts \(H\phi_K\).  With \(p=\sqrt{W}\phi_K\), it obtains
\(\langle p|H|p\rangle\) and \(\langle p|H|c\rangle\), then reconstructs
\(2\operatorname{Re}\langle p|H|q\rangle\) by subtraction and
\(\langle q|H|q\rangle\) from the complete-state energy identity.  Thus the
tail term is an algebraic reconstruction, not a separate \(Hq\) contraction.
Across both orbital bases and all 17 support sizes, the largest floating-point
discrepancy on recombining Eq.~(\ref{eq:s-decomposition}) is
\(6.1\times10^{-14}\,\Eh\).  The positive removed-coupling and negative
tail-quadratic contributions in Fig.~2 are
\(-2\operatorname{Re}\langle p|A|q\rangle/W\) and
\(-\langle q|A|q\rangle/W\), respectively.

\section{Hamiltonian-aware 512-determinant controls}

Five fixed-budget states separate inherited coefficients, coefficient
relaxation, orbital representation, and support choice.  In each orbital
representation, the top-512 amplitude support and inherited coefficients are
taken directly from the full Ritz tensor.  The script
\path{analysis/rediagonalize_topk_support.py} then builds the complete
512-by-512 Hamiltonian on exactly that support and diagonalizes it.  In the
benchmark representation, a separate selected-CI calculation constructs a
different support: it grows an approximately 80,000-determinant candidate
space by
Hamiltonian connections and Epstein--Nesbet estimates, ranks additions by
secular energy lowering, evaluates removals by the exact Cauchy secular
equation, accepts only energy-lowering swaps, and rediagonalizes every accepted
512-dimensional support.  Because the search uses seeded, finite-time
large-neighborhood restarts, this state is a reproducible candidate, not a
claim of a unique or globally optimal support.

For all three rediagonalized states, the dense subspace Hamiltonian is built by
\path{analysis/reference_slater_condon.py}, a compact implementation that
parses the FCIDUMP and applies explicit fermionic creation and annihilation
operators.  A separate PySCF embedding places the 512 coefficients into the
full \(15{,}504^2\) tensor and contracts \(Hc\).  Table~\ref{tab:k512} reports
the results.  The projected residual tests stationarity only within the listed
support; the full-space residual tests coupling to all excluded determinants.

\begin{table*}[h]
\caption{Fixed-budget 512-determinant controls.  Energy error is relative to
\(E_\star\).  ``Projected'' and ``full'' denote residual norms in the selected
subspace and complete active-space determinant space.  B and N denote the
benchmark and natural-orbital representations.}
\label{tab:k512}
\small
\begin{tabular}{lcccc}
\toprule
State & Energy (\(\Eh\)) & Error (m\(\Eh\)) & Projected residual (\(\Eh\)) & Full residual (\(\Eh\)) \\
\midrule
B: amplitude support, inherited coefficients & \(-116.000343445108\) & 605.266 & --- & 0.786 \\
B: same support, rediagonalized coefficients & \(-116.146308570453\) & 459.301 & \(2.42\times10^{-13}\) & 0.558 \\
B: Hamiltonian-guided support, rediagonalized & \(-116.370462675845\) & 235.146 & \(8.68\times10^{-14}\) & 0.215 \\
N: amplitude support, inherited coefficients & \(-116.436626144441\) & 168.983 & --- & 0.473 \\
N: same support, rediagonalized coefficients & \(\NaturalKFiveTwelveRelaxedEnergy\) & \(\NaturalKFiveTwelveRelaxedErrorMilli\) & \(\NaturalKFiveTwelveProjectedResidual\) & \(\NaturalKFiveTwelveFullResidual\) \\
\bottomrule
\end{tabular}
\end{table*}

The same-support energy from the explicit fermionic implementation and PySCF
differs by \(3.13\times10^{-13}\,\Eh\) in B and
\(\NaturalKFiveTwelveIndependentDifference\,\Eh\) in N; the corresponding
difference for the Hamiltonian-guided state is
\(1.71\times10^{-13}\,\Eh\).  The frozen wave-function CSVs have 512 unique
rows and squared norms equal to unity within \(4\times10^{-16}\).  The
construction and validation codes are distinct, and the reported energy is
always the independently contracted value.

\section{Software environment and reproduction map}

The production state used Python 3.12.1, NumPy 2.5.1, SciPy 1.18.0, h5py
3.14.0, and PySCF 2.14.0.  The PySCF and h5py wheel SHA-256 values are
\begin{quote}\small\ttfamily
37b0bccc55450311a55318cd643e851353331ddeab4fc0c0065e83c905e41502\\
0cbd41f4e3761f150aa5b662df991868ca533872c95467216f2bec5fcad84882.
\end{quote}
Thread controls \texttt{OMP\_NUM\_THREADS}, \texttt{MKL\_NUM\_THREADS},
\texttt{OPENBLAS\_NUM\_THREADS}, and \texttt{NUMEXPR\_NUM\_THREADS} were set
to 64.  The top-level \path{README.md} gives copy-paste commands for every
stage and records expected hashes at each boundary.

\begin{table*}[h]
\caption{Script-to-result map in the reproduction bundle.}
\label{tab:scripts}
\small
\begin{tabularx}{\textwidth}{@{}>{\raggedright\arraybackslash}p{0.34\textwidth}>{\raggedright\arraybackslash}p{0.23\textwidth}>{\raggedright\arraybackslash}X@{}}
\toprule
Script & Principal inputs & Principal outputs \\
\midrule
\path{code/initialize_fci_ritz.py} & Primary FCIDUMP & Stage-A vector and direct diagnostics \\
\path{code/continue_fci_checkpointed.py} & FCIDUMP, explicit NPY checkpoint & Atomic checkpoints, history JSONL, endpoint diagnostics \\
\path{code/verify_publication_state.py} & FCIDUMP, final vector & Direct/RDM energy, residual, norm, spin, occupations \\
\path{code/compare_hamiltonian_orbital_equivalence.py} & Two public FCIDUMPs & Orbital rotation and complete-tensor residuals \\
\path{analysis/analyze_state_compressibility.py} & FCIDUMP, representation-specific vector & Top-\(K\) CSV, decomposition CSV, occupations, summary JSON \\
\path{analysis/transform_ci_natural_basis.py} & Both FCIDUMPs, rotation, benchmark vector & Natural-basis vector and invariance JSON \\
\path{analysis/run_natural_basis_replay_v3.sh} & Natural FCIDUMP/vector, frozen occupations & Hash-locked two-basis replay and natural top-512 control \\
\path{analysis/rediagonalize_topk_support.py} & FCIDUMP, full vector, explicit reference code & Same-support ground state and two-path validation \\
\path{analysis/guided_512/exp_structure.py}, \path{opt3.py} & FCIDUMP and fixed search parameters & Hamiltonian-guided candidate pool and 512-state CSV \\
\path{analysis/validate_optimized_512.py} & FCIDUMP, guided 512 state, explicit reference code & Guided-support energy and residuals \\
\path{analysis/build_results_tex.py} & Gated compact JSON/CSV & Numerical macros used by both documents \\
\path{figures/build_figures.py} & Frozen compact CSV/JSON only & PDF, SVG, and 450-dpi PNG figures \\
\bottomrule
\end{tabularx}
\end{table*}

The compact reproduction bundle includes source code, job manifests,
environment locks, convergence history, verification JSON, the derived
natural-orbital FCIDUMP and rotation, all figure data, and SHA-256 manifests;
immutable locations and hashes for both public source FCIDUMPs are given above.
For anonymous scientific review the 1,922,992,256-byte final vector is supplied
as the separate artifact \path{fci_ritz_vector.npy} and identified by its full
hash above.  A persistent, unauthenticated public accession and final author
metadata must replace the review-only transfer before external submission.

\bibliography{references}

\end{document}
